% ============================================================
% Week 1 Lab Manual (LaTeX) — Concurrency & Distributed Systems
% Lebanese University, Faculty of Engineering
% Instructor: Dr. Mohamad Aoude
% ============================================================

\documentclass[11pt,a4paper]{article}

% -------------------------------
% Encoding & fonts (pdfLaTeX-safe)
% -------------------------------
\usepackage[utf8]{inputenc}
\usepackage[T1]{fontenc}
\usepackage{lmodern}

% -------------------------------
% Layout
% -------------------------------
\usepackage[a4paper,margin=1.6cm]{geometry}
\usepackage{parskip}

% -------------------------------
% Lists, tables, links
% -------------------------------
\usepackage{enumitem}
\usepackage{array}
\usepackage{booktabs}
\usepackage{longtable}
\usepackage{hyperref}
\usepackage{amssymb}
\usepackage{amsmath}

% -------------------------------
% Code
% -------------------------------
\usepackage{listings}
\usepackage{xcolor}

\lstset{
	basicstyle=\ttfamily\small,
	breaklines=true,
	columns=fullflexible,
	frame=single,
	rulecolor=\color{black},
	showstringspaces=false,
	tabsize=2
}

% -------------------------------
% Simple helpers
% -------------------------------
\newcommand{\course}{Concurrency \& Distributed Systems}
\newcommand{\univ}{Lebanese University}
\newcommand{\faculty}{Faculty of Engineering}
\newcommand{\instructor}{Dr.\ Mohamad Aoude}
\newcommand{\wk}{Week 1}
\newcommand{\labtitle}{The Waiting Problem}
\newcommand{\checkbox}{\(\square\)}
\newcommand{\checkedbox}{\(\blacksquare\)}


\usepackage{fancyhdr}
\usepackage{lastpage}
\usepackage{graphicx}
\usepackage{eso-pic}
\usepackage{qrcode}

% -------------------------------
% Header & Footer
% -------------------------------
\pagestyle{fancy}
\fancyhf{}

\fancyhead[L]{\textbf{ULFG --- SEM VIII}}
\fancyhead[R]{\textbf{Dr M AOUDE}}

\fancyfoot[L]{\textcopyright\ \the\year\ M AOUDE -- All Rights Reserved}
\fancyfoot[C]{Page \thepage\ /\ \pageref{LastPage}}
\fancyfoot[R]{\today}

\renewcommand{\headrulewidth}{0.4pt}
\renewcommand{\footrulewidth}{0.4pt}

% -------------------------------
% Synapse AI Watermark (pdfLaTeX-safe)
% Put synapse_ai_logo.png in same folder as .tex
% -------------------------------
\newcommand\SynapseWatermark{%
	\AddToShipoutPictureBG*{%
		\AtPageCenter{%
			\makebox(-400,550){%
				\color{black!6}%
				\includegraphics[width=0.25\paperwidth]{synapse_ai_logo_t.png}%
			}%
		}%
	}%
}

\begin{document}
	

\SynapseWatermark
	

\newpage	
		\begin{center}
		{\LARGE \textbf{\wk Environment Setup Manual: \labtitle}}\\[1mm]
		{\large \textit{\course} --- \univ, \faculty}\\[1mm]
		\textbf{Instructor:} \instructor \quad | \quad \textbf{Language:} Java (JDK 17+)
	\end{center}
	\vspace{6pt}
\vspace{6pt}
\begin{center}
	\begin{tabular}{cc}
		\qrcode[height=2.2cm]{https://github.com/maoude/concurrency-course}
		&
		\qrcode[height=2.2cm]{https://www.linkedin.com/in/mohamad-aoude-4aba1270/}
		\\[4pt]
		\small GitHub Repository
		&
		\small LinkedIn Profile
	\end{tabular}
\end{center}

	
	% ============================================================
	\section*{ — Week 1 Lab}
	
	This section provides a complete PowerShell-first setup guide for configuring
	the Week 1 Lab environment on Windows.
	
	The environment includes:
	\begin{itemize}
		\item Temurin JDK 25 (LTS)
		\item Visual Studio Code
		\item Java Extension Pack
	\end{itemize}
	
	\vspace{6pt}
	\hrule
	\vspace{10pt}
	
	% ============================================================
	\subsection*{Important — Administrative Privileges Required}
	
	\textbf{All installation commands must be executed in PowerShell running as Administrator.}
	
	\textbf{How to open PowerShell as Administrator:}
	
	\begin{enumerate}
		\item Open the Start Menu
		\item Type \texttt{PowerShell}
		\item Right-click on \textbf{Windows PowerShell}
		\item Select \textbf{Run as administrator}
	\end{enumerate}
	
	If PowerShell is not running as Administrator, system-wide PATH and JDK installation will fail.
	
	\vspace{6pt}
	\hrule
	\vspace{10pt}
	
	% ============================================================
	\subsection*{Step 0 — Verify or Install Winget}
	
	\textbf{0.1 Check if winget is installed}
	
	\begin{lstlisting}
		winget --version
	\end{lstlisting}
	
	If a version number is displayed, proceed to Step 1.
	
	\vspace{4pt}
	
	\textbf{0.2 If winget is NOT installed}
	
	Winget is part of the \textbf{App Installer} package in Windows 10/11.
	
	\textbf{Method 1 — Install via Microsoft Store:}
	
	\begin{itemize}
		\item Open Microsoft Store
		\item Search for \textbf{App Installer}
		\item Install or Update it
	\end{itemize}
	
	\vspace{4pt}
	
	\textbf{Method 2 — Install manually (offline method)}
	
	Download the latest App Installer package from:
	
	\begin{lstlisting}
		https://aka.ms/getwinget
	\end{lstlisting}
	
	After installation, close and reopen PowerShell, then verify:
	
	\begin{lstlisting}
		winget --version
	\end{lstlisting}
	
	\vspace{6pt}
	\hrule
	\vspace{10pt}
	
	% ============================================================
	\subsection*{Step 1 — Install JDK 25 (Temurin LTS)}
	
	\textbf{1.1 Install Temurin JDK 25}
	
	\begin{lstlisting}
		winget install -e --id EclipseAdoptium.Temurin.25.JDK --source winget
	\end{lstlisting}
	
	\vspace{6pt}
	
	% ============================================================
	\subsection*{Step 2 — Configure JAVA\_HOME and PATH}
	
	\textbf{2.1 Locate the JDK installation directory}
	
	\begin{lstlisting}
		Get-ChildItem "C:\Program Files\Eclipse Adoptium\" -Directory |
		Select-Object Name, FullName
	\end{lstlisting}
	
	Select the directory beginning with \texttt{jdk-25...}
	
	\vspace{4pt}
	
	\textbf{2.2 Set JAVA\_HOME and update system PATH}
	
	\begin{lstlisting}
		$jdk = "C:\Program Files\Eclipse Adoptium\jdk-25.0.x-hotspot"
		
		[Environment]::SetEnvironmentVariable("JAVA_HOME", $jdk, "Machine")
		
		$machinePath = [Environment]::GetEnvironmentVariable("Path","Machine")
		$newMachinePath = "$jdk\bin;$machinePath"
		[Environment]::SetEnvironmentVariable("Path", $newMachinePath, "Machine")
	\end{lstlisting}
	
	\vspace{4pt}
	
	\textbf{2.3 Refresh environment variables (no reboot required)}
	
	\begin{lstlisting}
		$env:Path = [Environment]::GetEnvironmentVariable("Path","Machine") +
		";" +
		[Environment]::GetEnvironmentVariable("Path","User")
		
		$env:JAVA_HOME =
		[Environment]::GetEnvironmentVariable("JAVA_HOME","Machine")
	\end{lstlisting}
	
	\vspace{4pt}
	
	\textbf{2.4 Verify installation}
	
	\begin{lstlisting}
		where.exe java
		where.exe javac
		java -version
		javac -version
		echo $env:JAVA_HOME
	\end{lstlisting}
	
	The first path shown must reference the Eclipse Adoptium directory.
	
	\vspace{6pt}
	\hrule
	\vspace{10pt}
	
	% ============================================================
	\subsection*{Step 3 — Install Visual Studio Code}
	
	\begin{lstlisting}
		winget install -e --id Microsoft.VisualStudioCode --source winget
	\end{lstlisting}
	
	Verify installation:
	
	\begin{lstlisting}
		code --version
	\end{lstlisting}
	
	If the command is not recognized, restart PowerShell.
	
	\vspace{6pt}
	\hrule
	\vspace{10pt}
	
	% ============================================================
	\subsection*{Step 4 — Install Java Extension Pack}
	
	\begin{lstlisting}
		code --install-extension vscjava.vscode-java-pack
	\end{lstlisting}
	
	Optional explicit installations:
	
	\begin{lstlisting}
		code --install-extension redhat.java
		code --install-extension vscjava.vscode-maven
		code --install-extension vscjava.vscode-gradle
	\end{lstlisting}
	
	\vspace{6pt}
	\hrule
	\vspace{10pt}
	
	% ============================================================
	\subsection*{Step 5 — Open and Run Week 1 Lab}
	
	\textbf{5.1 Open the project}
	
	\begin{lstlisting}
		cd D:\courses\concerent\week1-lab
		code .
	\end{lstlisting}
	
	\vspace{4pt}
	
	\textbf{5.2 Compile and run manually}
	
	\begin{lstlisting}
		cd D:\courses\concerent\week1-lab\lab1
		javac *.java
		java SingleThreadedServer
	\end{lstlisting}
	
	In a second terminal:
	
	\begin{lstlisting}
		java LoadTestClient
	\end{lstlisting}
	
	\vspace{6pt}
	\hrule
	\vspace{10pt}
	
	% ============================================================
	\subsection*{Step 6 — Sanity Test (Hello World)}
	
	\begin{lstlisting}
		cd D:\courses\concerent\week1-lab\lab1
		
		@'
		public class Hello {
			public static void main(String[] args) {
				System.out.println("Hello Java 25!");
			}
		}
		'@ | Set-Content .\Hello.java
		
		javac Hello.java
		java Hello
	\end{lstlisting}
	
	Successful output confirms correct configuration.
	
	\vspace{6pt}
	\hrule
	\vspace{6pt}
	
	\begin{center}
		\textit{A correct environment setup eliminates 90\% of concurrency debugging problems.}
	\end{center}

\end{document}
